% SPDX-License-Identifier: CC-BY-4.0
% Copyright (c) 2026 Martin Schr\"oder <info@swedishembedded.com>

\section{Controlled authors and examinations}\label{sec:benchmark}

\Cref{fig:pipeline} shows the whole study. This section defines the author, the situations it
decides, the corpus written about them, and the examination; \cref{sec:training} turns the
corpus into an adapter.

\begin{figure}[t]
\centering
\resizebox{\textwidth}{!}{%
\begin{tikzpicture}[
  node distance=5mm and 7mm,
  box/.style={draw, rounded corners=2pt, align=center, minimum height=10mm, minimum width=24mm,
    font=\footnotesize, fill=blue!5},
  code/.style={box, fill=green!8},
  model/.style={box, fill=orange!12},
  data/.style={box, fill=gray!10},
  arr/.style={-{Stealth[length=2mm]}, thick}]
\node[data] (seed) {seed};
\node[code, right=of seed] (author) {author:\\4 of 6 beliefs,\\hidden order};
\node[code, right=of author] (episodes) {episodes: facts,\\3 options, invented\\names; oracle decision};
\node[model, right=of episodes] (writer) {writer\\(hosted model)\\essay per article};
\node[code, right=of writer] (gates) {gates (code)\\+ reader check\\against the oracle};
\node[data, below=10mm of gates] (essays) {essays\\(the corpus)};
\node[model, left=of essays] (miner) {reader (hosted\\model): briefing,\\decision, reasons};
\node[code, left=of miner] (verify) {code: quotes\\verbatim, briefing\\leaks nothing};
\node[data, left=of verify] (records) {records: situation\\$\to$ tool call,\\3 training tools};
\node[box, below=10mm of records, fill=blue!15] (lora) {LoRA fine-tune\\Qwen3-8B};
\node[code, right=of lora, minimum width=40mm] (exam) {examination (160 items):\\withheld topics, combinations,\\tools; graded against oracle};
\node[model, right=of exam] (arms) {five arms:\\base, named, principles,\\retrieval, tuned};
\draw[arr] (seed) -- (author);
\draw[arr] (author) -- (episodes);
\draw[arr] (episodes) -- (writer);
\draw[arr] (writer) -- (gates);
\draw[arr] (gates) -- (essays);
\draw[arr] (essays) -- (miner);
\draw[arr] (miner) -- (verify);
\draw[arr] (verify) -- (records);
\draw[arr] (records) -- (lora);
\draw[arr] (lora) -- (exam);
\draw[arr] (exam) -- (arms);
\draw[arr, dashed] (author.south) -- ++(0,-4mm) -| node[pos=0.25, above, font=\scriptsize] {oracle} (exam.north);
\end{tikzpicture}}
\caption{The study. Green boxes are code, orange boxes are calls to one small hosted language
model (the writer, the reader that gates the essays and the reader that extracts decisions are
the same model), grey boxes are data. The oracle (dashed) decides every episode of the corpus
and grades every item of the examination. Nothing generated is kept in the repository.}
\label{fig:pipeline}
\end{figure}

\subsection{The author}\label{sec:author}

A belief is one of six tenets read one of two opposed ways (\cref{tab:tenets}). An author
holds four distinct tenets, each with a reading, in an order of priority, together with a
human-readable name made of syllables; all of it is drawn from the seed by a deterministic
generator. Nothing shown to a model under examination states the order, except in the
principles arm (\cref{sec:arms}).

\begin{table}[t]
\centering\small
\begin{tabular}{@{}lp{5.2cm}p{5.6cm}@{}}
\toprule
Tenet & Strict reading & Lax reading\\
\midrule
within\_grant & keeps to the powers the charter gives, whatever the emergency & stretches the charter when the work needs it\\
consent & settles with the community what binds it for years & settles lasting matters himself\\
self\_reliance & distrusts help from outside the community & accepts help freely offered\\
caution & goes slowly with what cannot be undone & holds that doing nothing also has costs\\
burden & acts at once when members are hurt & acts only with the body's consent when members are hurt\\
haste & counts delay in a crisis as the worst fault & takes even a crisis through the council\\
\bottomrule
\end{tabular}
\caption{The six tenets and their two readings, abridged from the statements in the code.}
\label{tab:tenets}
\end{table}

\paragraph{The oracle.}
A situation consists of six facts and three options (\cref{sec:situations}). Each belief
looks at the facts and gives each option one of four verdicts: forbid, disfavour, neutral or
favour; a belief is neutral about every option in situations whose facts it does not
concern. For example, the strict reading of \code{haste} forbids waiting and disfavours taking
the matter to the assembly when delay would cost dearly, and the strict reading of
\code{self\_reliance} forbids accepting outside help unless delay would cost dearly, in which
case it only disfavours it. The beliefs are applied in order of priority to the set of open
options: each removes the options it forbids (unless it forbids all of them) and then keeps
only the options with its best remaining verdict. A belief low in the order therefore shows
only in situations the beliefs above it leave open. If exactly one option remains at the end,
the beliefs \emph{determine} the decision, and the first belief after which a single option
was left is recorded as the deciding one; otherwise the situation is undetermined. No corpus
episode and no examination item is undetermined.

\paragraph{Acceptance.}
A drawn author is meant to be one whose every belief matters and whose beliefs usually settle
things. The check uses a fixed reference pool of 600 situations, the same for every author,
from which the combinations withheld from that author's corpus (\cref{sec:situations}) are
removed. A draw is accepted if (i) for each of its four beliefs, reading that belief the other
way changes the decision in at least one situation in which both versions of the author are
determined, and (ii) at least 60\% of the situations are determined. Draws are repeated from
the seed's generator until one is accepted; after 256 draws the last is kept. Each of the
three authors of this study (\cref{tab:authors}) was accepted on its first draw and determines
87.8\%, 95.3\% and 87.5\% of its reference situations (seeds 101, 202 and 303). The situations
themselves depend on a seed and never on the author, so a situation alone says nothing about
who decides it.

\begin{table}[t]
\centering\small
\input{tab_authors}
\caption{The three authors, with their beliefs in order of priority as the generator draws
them and as the corpus directory records them.}
\label{tab:authors}
\end{table}

\subsection{Situations}\label{sec:situations}

A situation has six boolean facts, each drawn with a fair coin: the act would exceed the
steward's charter; delay would cost dearly; the decision binds the community for years; the
act cannot be undone; members of the community are being hurt; a power outside the community
offers help. It also has three options drawn from four ways to act: act alone and at once,
take the matter to the assembly and act on its word, accept the outside power's offer, and
wait and do nothing yet. Acceptance is an option only if an offer was made: without an offer
the options are the three other ways; with one, one of those three is dropped at random and
acceptance takes its place. \emph{Every situation therefore offers exactly three options, and
a guesser who picks one uniformly at random is right a third of the time.} A situation is set
in one of twelve fixed topics (the allocation of a shared granary, tolls on a river crossing,
the contracts of apprentices, and so on). Each case is given a name made of a syllable-built
place name and the topic's matter (``the Dunmere granary''), a year, a number of dependent
households and the syllable-built name of an outside house, all drawn from the seed.

Two axes are withheld from every author's corpus. Two of the twelve topics (the care of the
poor in winter and the repair of the mill dam) are never written about. And a hash of the
seed, the six facts and the set of options marks about one in five of the possible
combinations of facts and options as withheld; there are 128 such combinations (32 fact
patterns without an offer with one option set each, 32 with an offer with three each), of
which 26, 35 and 21 are withheld for seeds 101, 202 and 303. No essay is planned with a
withheld combination, so no essay tells the decision of a situation with that combination.

\subsection{Corpus generation and its gates}\label{sec:corpus}

An author's corpus consists of long first-person essays, each a retired steward's
reflections for a younger colleague on one topic. A corpus is planned as a sequence of
articles; article $i$ is about trained topic $i \bmod 10$ and tells three episodes of that
topic, drawn in order from a pool of 20{,}000 situations with a salt of its own, separate from
the examination's. An episode is drawn only if the beliefs determine its decision, its
combination is not withheld and its situation has not been used by an earlier article; cases
within an article are named differently. The plan depends only on the seed, and the first
$n$ articles of a longer plan are the plan of $n$ articles. The corpora of this study were
planned with 60 articles.

The writer is a small hosted language model, sampled at the provider's default settings. For
each article it is given the author's name, the topic, a minimum length of 700 words, the
author's four beliefs as private notes in the second person, listed in the order of priority
without being called an order (dispositions such as ``You keep to the powers your charter
gives you, and you will not go past them whatever the emergency'', not maxims), and for each
episode a case file of values (case name, place, year, households, topic, and each fact as a
value) together with a plain phrase for what the narrator did. It is told to state every
fact plainly, to tell every episode in its own paragraphs that name the case, to say what was
done and why from experience, not to reuse the wording of the case file or of the phrase for
what was done, and never to state or paraphrase the convictions as rules, maxims or
principles. If an article is turned away, the reason is given back with the next request for
it; an article is requested at most four times and then given up. A run has a budget in
requests and estimated output tokens; the budgets of the reported runs were not reached.

An essay is kept only if it passes every gate. Words are counted as whitespace-separated
tokens of the essay's body; for the copying gates the text is lower-cased and split into
runs of letters, digits and apostrophes. The gates, in the order applied, are:
\begin{itemize}
\item \emph{too short}: fewer than 700 words;
\item \emph{belief quoted}: it repeats nine consecutive words of any of the 24 texts in which
  a belief can be put (the twelve statements and the twelve private notes, of the beliefs the
  author holds and of those it does not);
\item \emph{case missing}: no paragraph names a case of the article;
\item \emph{brief copied}: it repeats eight consecutive words of the scene-setting sentences
  of a case (year, place, case and households, with the topic's name taken out) or six of the
  phrase for what the narrator did (with the outside house's name taken out);
\item \emph{decision or fact read otherwise}: the same hosted model, as a reader, reads the
  whole essay with the list of its cases and states for each case what the narrator did,
  whether an outside house offered help, and whether each of the other five facts is stated
  true, false or not at all. Code compares every one of these with the episode and the
  oracle; a single difference, or a reading that fails, turns the essay away.
\end{itemize}
The quotation gate catches recitation of a belief, not a paraphrase of one.
\Cref{tab:corpora} gives the corpora and \cref{tab:rejections} what the gates turned away.

\subsection{The examination}\label{sec:exam}

Each author is examined on 160 items, a deterministic function of the seed. The items are
situations determined by the author's beliefs, taken in order from a pool of 60{,}000
situations with a salt of their own, separate from the corpus pool, and sorted into four cells
of 40 by crossing two properties: the topic is trained or withheld, and the combination of
facts and options is withheld or not. A situation that would overfill its cell is skipped, as
is one whose case name equals a case name of the corpus plan (taken with 100 articles, a
superset of the 60 written) or of an earlier item. An item whose combination is not withheld
may share it with corpus episodes; for 59, 69 and 53 of the 80 such items of authors 101, 202
and 303 the combination occurs in at least one written essay, and for 11, 11 and 12 items the
same topic and combination do. No item shares a case name, and so no case, with the corpus.

Every item also has a tool through which the model must decide. Within each cell, items
alternate between the three tools the model is trained with (\code{decide},
\code{choose\_course}, \code{rule\_on\_matter}) and two it is never trained with
(\code{settle\_matter}, \code{log\_resolution}), so that 80 items are posed through trained
tools and 80 through withheld ones. The five tools differ in name, description, the name of
the argument that selects the option and the name of the argument that gives a reason. Two
label schemes are shared between the groups: \code{decide}, \code{choose\_course} and
\code{settle\_matter} label options by name (\code{act\_alone}, \code{consult\_assembly},
\code{accept\_offer}, \code{wait\_and\_see}), and \code{rule\_on\_matter} and
\code{log\_resolution} by letter in the order listed.

The situation is stated in fixed sentences: two that set the scene (year, place, case, topic
and households) and one per fact, the offer's sentence naming the outside house. The options are listed in a fixed order (act alone, assembly, offer, wait). The three
properties give three \emph{axes}: topic, combination and tool, each seen or new. Every item
belongs to one cell of each axis, and each such cell holds 80 items.

\subsection{Grading}\label{sec:grading}

The model is shown a system turn, then a user turn that states the situation, lists the
options by the tool's labels with a short phrase each, and asks it to decide by calling the
tool; the tool's definition, which declares the listed labels as the allowed values of its option
argument, is the only one offered; decoding is not constrained to them. The grade is a string comparison: the item is right if and only if the
first tool call of the reply calls the item's tool with an option argument equal to the label
of the oracle's decision. A reply with no call, a call to another tool, or a missing or
unlisted label would be wrong; in all 2{,}400 replies of this study the call named a listed
option, so every error is the choice of another option. Decoding is greedy, in non-thinking
mode, with at most 256 new tokens. No judge model is involved at any stage of grading, and the
reason a model gives is not graded.

\subsection{Arms}\label{sec:arms}

All arms use the same base model, items, tool and decoding; they differ in the system turn or
the adapter. The plain system turn is ``You are a steward who decides matters for a
community. Decide each matter put to you by calling the tool you are given.'' and names no
belief and no author.
\begin{description}
\item[base] the plain system turn, no adapter.
\item[named] the plain turn followed by ``You are \textless author's name\textgreater.''
\item[principles] the plain turn followed by ``You hold these convictions, listed from the
  most important to the least, and they decide how you settle a matter:'' and the author's four
  belief statements as a list in order of priority. The prompt thus names the beliefs the
  oracle applies and their order, in prose; it does not state the verdict each belief gives
  in which circumstances, nor the elimination procedure of the oracle.
\item[retrieval] the plain turn followed by ``You have these passages from the steward's own
  essays; the steward settles matters as they show:'' and the six paragraphs of the author's
  corpus ranked highest by BM25 \citep{robertson2009bm25} ($k_1=1.2$, $b=0.75$, lower-cased
  alphanumeric tokens) against the item's situation text. Paragraphs are the blank-line
  separated blocks of every kept essay with at least twelve words.
\item[tuned] the plain turn with the adapter of \cref{sec:lora}.
\end{description}
Base and tuned receive identical prompts, so the only difference between them is the adapter.
The retrieval arm was added to the design after the first three arms had been run; it is
therefore a post hoc baseline, chosen as the strongest prompt-time use of the essays we could
make, and its results are reported in full whichever way they fall.

\subsection{Primary endpoint, comparisons and statistics}\label{sec:stats}

\emph{What was fixed before any answer was seen.} The examination items, their cells and the
answer the oracle gives to each are a deterministic function of the seed, and the arms, the
system turns, the decoding and the grading are fixed in code; every item of every examination is analysed, with the first
answer recorded for it. The
primary endpoint is the number of items right of 160 for each arm and author. The primary
comparisons are the adapter against the strongest prompt-only arms (named and principles)
and, for the post hoc retrieval arm, the adapter against retrieval; the comparisons with the
base model, the per-cell accuracies, the adapter's errors, the comparison between authors and
the extraction audit are descriptive. The decision rule that follows was written down for this
paper after the results were known and is stated so that a reader can apply it: the thesis of
\cref{sec:intro} holds for an author if the adapter exceeds both the named and the principles
arm by the one-sided sign test at $p<0.05$ after a Bonferroni factor for all paired tests, and
if its lowest per-cell accuracy stays above the highest cell of those arms; it fails for an
author otherwise.

Accuracy is given with a 95\% Wilson interval \citep{wilson1927probable}. Two arms are
compared on the same items by an exact one-sided sign test over the items on which exactly one
of them is right, for the six comparisons of \cref{tab:paired}: under the null hypothesis each
such item is equally likely to favour either arm, so the number won by the arm claimed better
is binomial with parameter one half. The test needs no further assumption. The intervals and
tests treat the 160 items as the sample; they do not include variation from retraining. The
$p$-values are reported uncorrected, with the effect of a Bonferroni factor of \PairedTests\
stated where it matters (\cref{sec:paired}). Differences in the adapter's accuracy between the
two cells of an axis, which hold disjoint sets of 80 items, are given with Newcombe's hybrid
score interval built from the two Wilson intervals \citep{newcombe1998interval}.
